% ID: SEL-D0082-OB-E03
% Tanda S082.
Sean $A$ y $B$ comprar en cada supermercado.

\textbf{a)} Por inclusión y exclusión,
\[
 P(A\cup B)=0{,}58+0{,}35-0{,}12=\boxed{0{,}81}.
\]
\textbf{b)} No comprar en ninguno es el complementario, de probabilidad $1-0{,}81=\boxed{0{,}19}$.

\textbf{c)} Comprar en exactamente uno comprende solo A o solo B. Su probabilidad es $(0{,}58-0{,}12)+(0{,}35-0{,}12)=\boxed{0{,}69}$.

\textbf{d)} La probabilidad conjunta de comprar en A y no en B es $0{,}46$, y $P(B^c)=0{,}65$. Por la definición de probabilidad condicionada,
\[
 P(A\mid B^c)=\frac{0{,}46}{0{,}65}=\boxed{\frac{46}{65}}\approx0{,}7077.
\]
